
\documentclass{article}
\usepackage{colortbl}
\usepackage{makecell}
\usepackage{multirow}
\usepackage{supertabular}

\begin{document}

\newcounter{utterance}

\centering \large Interaction Transcript for game `cladder', experiment `full\_v1.5\_default', episode 2340 with selfrefine{-}llama{-}3.1{-}8b{-}instruct.
\vspace{24pt}

{ \footnotesize  \setcounter{utterance}{1}
\setlength{\tabcolsep}{0pt}
\begin{supertabular}{c@{$\;$}|p{.15\linewidth}@{}p{.15\linewidth}p{.15\linewidth}p{.15\linewidth}p{.15\linewidth}p{.15\linewidth}}
    \# & \multicolumn{2}{c}{Player} && \multicolumn{2}{c}{Game Master} \\
    \hline

    \theutterance \stepcounter{utterance}  
    & & & \multicolumn{4}{p{0.6\linewidth}}{
        \cellcolor[rgb]{0.9,0.9,0.9}{
            \makecell[{{p{\linewidth}}}]{
                \texttt{\tiny{[P1$\langle$GM]}}
                \texttt{You are an expert in causal inference. The following question is not a typical commonsense query, but rather a meticulously designed question created by a professor specializing in causal inference, intended to assess the students' mastery of the course content.} \\
\texttt{Imagine a self{-}contained, hypothetical world with only the following conditions, and without any unmentioned factors or causal relationships: Eating citrus has a direct effect on vitmain C. Vitmain C has a direct effect on scurvy. For patients not consuming citrus, the probability of scurvy is 73\%. For patients consuming citrus, the probability of scurvy is 54\%. Does eating citrus positively affect scurvy through vitmain C?} \\
\texttt{Start your answer with 'Yes' or 'No', followed by additional reasoning or evidence to support your explanation.} \\
            }
        }
    }
     \\ \\

    \theutterance \stepcounter{utterance}  
    & & \multicolumn{4}{p{0.6\linewidth}}{
        \cellcolor[rgb]{0.95,0.95,0.95}{
            \makecell[{{p{\linewidth}}}]{
                \texttt{\tiny{[P1$\leftrightarrow$P]}}
                \texttt{{[}DRAFT{]}} \\
\texttt{No.} \\
\\ 
\texttt{To assess whether eating citrus positively affects scurvy through vitamin C, we need to consider the potential causal relationships and their directions. } \\
\\ 
\texttt{First, let's establish the given causal relationships: } \\
\\ 
\texttt{1. Eating citrus → Vitamin C (direct effect)} \\
\texttt{2. Vitamin C → Scurvy (direct effect)} \\
\\ 
\texttt{We are also given the probabilities of scurvy for patients not consuming citrus (73\%) and those consuming citrus (54\%). } \\
\\ 
\texttt{To determine whether eating citrus positively affects scurvy through vitamin C, we need to consider the potential causal relationship between eating citrus and scurvy, and whether it is mediated by vitamin C. } \\
\\ 
\texttt{One way to approach this is to use the backdoor criterion, which states that a path from the cause (eating citrus) to the effect (scurvy) is blocked if we can control for all the variables that lie on the path between the cause and the effect. } \\
\\ 
\texttt{In this case, the path from eating citrus to scurvy is: eating citrus → vitamin C → scurvy. } \\
\\ 
\texttt{To block this path, we need to control for vitamin C. However, we are given the probabilities of scurvy for patients not consuming citrus (73\%) and those consuming citrus (54\%). } \\
\\ 
\texttt{We can use the formula for the backdoor criterion to determine whether eating citrus positively affects scurvy through vitamin C:} \\
\\ 
\texttt{P(scurvy | eating citrus) = P(scurvy | vitamin C) * P(vitamin C | eating citrus) + P(scurvy | not eating citrus) * P(not vitamin C | eating citrus)} \\
\\ 
\texttt{We know that P(scurvy | not eating citrus) = 0.73 and P(scurvy | eating citrus) = 0.54. } \\
\\ 
\texttt{However, we do not have the direct probabilities of vitamin C given eating citrus, P(vitamin C | eating citrus), or not eating citrus, P(not vitamin C | eating citrus). } \\
\\ 
\texttt{But we can use the fact that the probability of scurvy given eating citrus is 54\% to infer the probability of vitamin C given eating citrus. } \\
\\ 
\texttt{Since eating citrus has a direct effect on vitamin C, we can assume that the probability of vitamin C given eating citrus is close to 1 (i.e., P(vitamin C | eating citrus) ≈ 1). } \\
\\ 
\texttt{Similarly, we can assume that the probability of not vitamin C given eating citrus is close to 0 (i.e., P(not vitamin C | eating citrus) ≈ 0).} \\
\\ 
\texttt{Substituting these values into the formula, we get:} \\
\\ 
\texttt{P(scurvy | eating citrus) ≈ P(scurvy | vitamin C) * 1 + 0.73 * 0} \\
\\ 
\texttt{Since P(scurvy | eating citrus) = 0.54, we can solve for P(scurvy | vitamin C):} \\
\\ 
\texttt{0.54 ≈ P(scurvy | vitamin C)} \\
\\ 
\texttt{However, we are given that vitamin C has a direct effect on scurvy} \\
            }
        }
    }
    & \\ \\

    \theutterance \stepcounter{utterance}  
    & & \multicolumn{4}{p{0.6\linewidth}}{
        \cellcolor[rgb]{0.95,0.95,0.95}{
            \makecell[{{p{\linewidth}}}]{
                \texttt{\tiny{[P1$\leftrightarrow$P]}}
                \texttt{{[}SELF{-}FEEDBACK{]}} \\
\texttt{FORMAT: ok} \\
\texttt{RULES: ok} \\
\texttt{STRATEGY: problem} \\
\texttt{REASON: The answer incorrectly concludes that eating citrus does not positively affect scurvy through vitamin C, despite the provided calculations suggesting a positive effect.} \\
            }
        }
    }
    & \\ \\

    \theutterance \stepcounter{utterance}  
    & \multicolumn{4}{p{0.6\linewidth}}{
        \cellcolor[rgb]{0.561,0.737,0.561}{
            \makecell[{{p{\linewidth}}}]{
                \texttt{\tiny{[P1$\rangle$GM]}}
                \texttt{No.} \\
\\ 
\texttt{The backdoor criterion is not met because we cannot control for vitamin C. The formula for the backdoor criterion is:} \\
\\ 
\texttt{P(scurvy | eating citrus) = P(scurvy | vitamin C) * P(vitamin C | eating citrus) + P(scurvy | not eating citrus) * P(not vitamin C | eating citrus)} \\
\\ 
\texttt{However, we assumed that P(vitamin C | eating citrus) ≈ 1 and P(not vitamin C | eating citrus) ≈ 0. This means that the formula simplifies to:} \\
\\ 
\texttt{P(scurvy | eating citrus) ≈ P(scurvy | vitamin C)} \\
\\ 
\texttt{Since P(scurvy | eating citrus) = 0.54 and P(scurvy | not eating citrus) = 0.73, we can conclude that P(scurvy | vitamin C) < 0.73. This implies that eating citrus has a negative effect on scurvy through vitamin C, not a positive one.} \\
            }
        }
    }
    & & \\ \\

    \theutterance \stepcounter{utterance}  
    & & & \multicolumn{2}{p{0.3\linewidth}}{
        \cellcolor[rgb]{0.9,0.9,0.9}{
            \makecell[{{p{\linewidth}}}]{
                \texttt{\tiny{[GM$|$GM]}}
                \texttt{No} \\
            }
        }
    }
    & & \\ \\

    \theutterance \stepcounter{utterance}  
    & & & \multicolumn{2}{p{0.3\linewidth}}{
        \cellcolor[rgb]{0.9,0.9,0.9}{
            \makecell[{{p{\linewidth}}}]{
                \texttt{\tiny{[GM$|$GM]}}
                \texttt{no} \\
            }
        }
    }
    & & \\ \\

    \theutterance \stepcounter{utterance}  
    & & & \multicolumn{2}{p{0.3\linewidth}}{
        \cellcolor[rgb]{0.9,0.9,0.9}{
            \makecell[{{p{\linewidth}}}]{
                \texttt{\tiny{[GM$|$GM]}}
                \texttt{game\_result = WIN} \\
            }
        }
    }
    & & \\ \\

\end{supertabular}
}

\end{document}
